% Symbolic validation component of Evaluation.
\label{sec:model-validation}

이 검사 프로그램은 거래에 들어온 기록과 새로 나온 기록의
관계, 기록의 소비 여부와 동결 상태만 Python 3의 자료구조로
표현한다. 암호문과 소비 식별자는 실제 암호 연산 대신
서로 구별되는 이름표로 나타낸다. 영지식증명·암호화·인증이나
원장 거래는 실행하지 않는다.

추적 결과를 확인하기 위해 같은 이력을 두 방식으로 계산했다.
첫째는 제안한 감사 절차처럼 감사 자료와 소비 여부를 조회해
연결된 기록을 찾는다. 둘째는 거래마다 적어 둔 입력·출력
목록만으로 연결 관계를 새로 만든 뒤, 도달할 수 있는 기록을
모두 찾는다. 둘째 계산은 첫째 계산의 결과를 사용하지 않는다.
분기·합류·발급·종료, 이전 회수, 질량 0 출력 등을 포함한
네 개의 정해진 거래 이력을 사용했다. 각 이력에서 거래가
하나씩 추가되는 모든 시점마다, 그때 존재하는 기록과 거래를
각각 출발점으로 삼아 앞뒤 추적 결과를 비교했다.
기록·연결·아직 소비되지 않은 대상에 대한 비교 334건에서
두 방식의 결과가 일치했다~\cite{zkdpp-symbolic-model-2026}.
하지만 두 방식이 같이 틀리면 이 비교로 발견할 수 없다.
실제로 Entry 거래 자체에서 후속 추적을 시작하면 둘 다
Entry를 결과에 넣는다. 명세에서는 Entry의 출력부터
포함해야 한다. 따라서 334건의 일치가 명세 전체와의
일치를 뜻하지는 않는다.

그림~\ref{fig:worked-flow}의 원료 $A$에서 추적을 시작하면
아직 소비되지 않아 동결해야 할 대상은 $Z,E,W$다.
합류 거래에 함께 들어온 다른 원료 $X$의 과거로는
거슬러 올라가지 않는다. 이미 발급된 $K$도 더 소비할
수 없으므로 동결 대상이 아니다. 반면 $Z$는 질량이 0이고
$W$는 폐기물이지만 둘 다 아직 소비할 수 있어 포함된다.
$V$가 인수되기 전 시점에서는 대상이 $C,V,Z$다.
감사를 시작한 뒤 $V$가 먼저 인수되더라도, 처음 정한
대상 목록에서 $V$를 빼지 않는다. 동결을 마치지 못한 것으로
보고한다.

전체 검사 408건에는 앞의 추적 비교 334건과
동결 상태·오류 처리·회계 규칙 검사가 포함된다.
정해진 사례에서 동결되거나 사용 금지된 기록을 거래 입력으로
쓰면 거부했고, 이전 증표의 인수와 회수 중 먼저 수락된
한쪽만 그 증표를 소비할 수 있었다. 감사 자료가 빠지거나
형식이 틀리면 추적이 실패했다. 동결 대상이 빠졌거나
일부만 동결된 경우, 제출 결과가 불분명하거나 동결 전에
대상이 소비된 경우에는 완료를 보고하지 않았다.
각 대상의 동결 영수증이 따로 존재해도 모든 대상이
\emph{같은 시점에} 차단된 적이 없으면 완료로 보지 않았다.
질량 0 가공 출력에 탄소 잔여를 주려는 이전 사례는 거부하고,
잔여가 양의 질량 출력에 배분되는지 확인했다. 서로 다른 소유
비밀이 같은 소유자를 나타내는 사례는 식별자 가정의 한계를
보여 주었다.

이 검사는 정해진 몇 가지 사례에서 규칙의 동작을 살펴본
것이다. 실제 암호·원장 구현이 명세를 따른다는 증거도,
모든 가능한 거래 이력에서 올바르게 작동한다는 증거도 아니다.
